\section{2026-08-11 — El apartado legal se ubica en el capítulo de metodología, no en el de introducción}

El mapa de correspondencias entre el wiki y el documento asigna el contenido de \texttt{wiki/proyecto/} al Capítulo~1. Siguiendo esa regla, el análisis de restricciones legales y éticas debería haberse incorporado a la introducción. Se resolvió apartarse de la regla y ubicarlo como sección del capítulo de metodología de desarrollo, junto a las decisiones de diseño que fundamenta.

La razón es que ese análisis dejó de ser un contexto normativo para convertirse en la justificación de decisiones concretas de diseño. La conservación sin plazo se apoya en el artículo 4 inciso 7 de la Ley 25.326 y determina que ninguna entidad del modelo de datos incorpore un atributo de fecha de purga. El alcance de la supresión que fija el artículo 16 determina qué tres atributos se eliminan y por qué la dirección de la publicación no se persiste. La exención de consentimiento del artículo 11 inciso 3 para los supuestos de disociación es la que convierte a la anonimización de la entrega hacia organizaciones en un requerimiento de prioridad imprescindible en lugar de una declaración de intenciones. Y la eximente de atribución fiel del artículo 113 del Código Penal determina la redacción del nivel más severo del indicador y el rótulo del segundo módulo del sistema.

Ubicado en la introducción, el lector encuentra un análisis normativo veinte páginas antes que las decisiones que ese análisis explica, y ninguno de los dos textos se sostiene solo. Ubicado en el capítulo de metodología, cada apartado remite a la sección de modelo de datos o a la de tecnologías que lo materializa, y esas referencias cruzadas funcionan.

La regla del mapa se conserva para el resto del contenido de \texttt{wiki/proyecto/}. La excepción alcanza únicamente a este archivo y queda registrada acá para que la próxima revisión no la interprete como un descuido.

\bigskip

\noindent\textbf{Nota sobre la numeración de capítulos.} Mientras el Capítulo~3 permanezca comentado en \texttt{main.tex}, la metodología de desarrollo se imprime numerada como Capítulo~3 y no como Capítulo~4. Es el comportamiento esperado y se corrige solo al incorporar el capítulo de descripción. Ninguna referencia cruzada del documento nombra el número de ese capítulo, de modo que la numeración provisional no produce referencias sin resolver.
